\section{Intensive composition lemmas}
\label{sec:lemmas}

The deployment-safety theorem of \S\ref{sec:main_theorem} composes
nine lemmas: four general composition lemmas (1--4) and a five-part
family covering substrate-targeting adversarial events (5a--5e).
This section states and proves each.

The lemma dependency structure is:
\begin{itemize}
\item Lemma~1 establishes that intensive bounds compose to intensive
  bounds under co-evolution; this is used by every subsequent lemma.
\item Lemmas~2 and~3 bridge the inherited source-paper apparatus
  (\citep[Phase~Redundancy]{lasser2026phase} Lyapunov, \citep[Need~Sufficiency]{lasser2026paper8b} HHI) to \citep[Microfoundation]{lasser2026micro}'s static Goodhart
  bound.
\item Lemma~4 lifts \citep[Exogenous~Verification]{lasser2026exo}'s SPRT mean-time bound to a tail-bounded
  lead-time guarantee.
\item Lemma~5a establishes a substrate floor for $\rext$ given
  $\mindep \geq m^*$.
\item Lemma~5b restricts the SPRT detection class to four-channel
  observable adversarial strategies and supplies channel-specific
  KL floors.
\item Lemma~5c is the static safe region from \citep[Horizon~Aware]{lasser2026horizon}'s minimax-
  form lift, scoped to the cooperative-overlap regime.
\item Lemma~5d composes Lemma~5b's KL floor with Lemma~4's tail
  bound to give the operational lead-time guarantee.
\item Lemma~5e extends Lemma~5b's machinery to environment-side
  observables under invariant $I_8$.
\end{itemize}

\subsection{Lemma 1 (Intensive composition under co-evolution)}
\label{sec:lem1}

\begin{definition}[Intensive in $|P|$ over deployment class $\mathcal{D}$]
\label{def:intensive_paper10}
Let $\mathcal{D}$ be a deployment class characterized by fixed
operational parameters (threat model, training-discipline regime,
verification infrastructure, governance protocol).  The operational
parameters are fixed before $|P|$ is varied; they may not include
arbitrary state-dependent quantities chosen to absorb the bound.

A bound $B$ depending on the deployment state is \emph{intensive in
$|P|$ over $\mathcal{D}$} if there exists a constant $C \geq 0$
that may depend on $\mathcal{D}$'s operational parameters but is
independent of $|P|$, such that $B \leq C$ for all valid deployment
states in $\mathcal{D}$.
\end{definition}

\begin{lemma}[Intensive composition under co-evolution]
\label{lem:1}
Let $X$ be a non-residual gap quantity (e.g., $\epsnonresComposed$,
or $f(\epsilon_{\mathrm{safe}})$ after Lemma~\ref{lem:2}'s
substitution).  Under conditions (C6) (bounded-Lipschitz alignment
property), (C7) (bounded co-evolution), and (C8) (per-capability
admissibility) of Theorem~\ref{thm:deployment_safety}, the
composition rule
\[
  B \;=\; \mathrm{Lip}(g) \cdot \bigl(X + \lambda \cdot \epsfloor\bigr),
  \qquad \lambda \in [0, 1]
\]
is intensive in $|P|$ over the deployment class $\mathcal{D}$:
there exists a constant $C^* \geq 0$ depending on $\mathcal{D}$'s
operational parameters but independent of $|P|$, such that
$B \leq C^*$ for all valid deployment states in $\mathcal{D}$.
\end{lemma}

\begin{proof}[Proof outline; full proof in
Appendix~\ref{app:proof_lem1}]
Under (C8) per-capability admissibility, each capability's gap is
allocated a ceiling with disjoint per-channel sub-shares; the
sup-norm $\epsnonresChannel{j}$ over capability-space inherits the
sub-share ceiling $K_j$, which is $|P|$-independent over
$\mathcal{D}$. The four-channel decomposition under co-evolution
gives $\epsnonresComposed \leq \sum_{j=1}^{4} \epsnonresChannel{j}
+ (\epsnonresCoev)_+$, where the positive-part co-evolution term is
bounded by (C7) bounded co-evolution
(Assumption~\ref{ass:bounded_coev}). The residual floor
$\epsfloor = \volR(\ResS)/\volR(\Pproxy) \leq 1$. Multiplying by
$\mathrm{Lip}(g) \leq K_{\mathrm{Lip}}$ from (C6) preserves
intensivity.  The composed constant is $C^* = K_{\mathrm{Lip}}
\cdot (K_X + K_{\mathrm{floor}})$.
\end{proof}

\textbf{Connection to Theorem~\ref{thm:deployment_safety}'s Layer~1.}
Lemma~\ref{lem:1} establishes that the composed slack term
$X + \lambda \cdot \epsfloor$ is intensive in $|P|$ over
$\mathcal{D}$, and that multiplying by $\mathrm{Lip}(g)$ preserves
intensivity.  Theorem~\ref{thm:deployment_safety}'s Layer~1 invokes
this with $X = f(\epsilon_{\mathrm{safe}})$ after Lemma~\ref{lem:2}'s
substitution; the resulting $h_{\mathrm{static}}(\theta) =
f(\epsilon_{\mathrm{safe}}) + \lambda \cdot \epsfloor$ is intensive
by Lemma~\ref{lem:1}, delivering the capability-magnitude-
independence property of the deployment claim.

\subsection{Lemma 2 (Lyapunov-to-Goodhart bridge)}
\label{sec:lem2}

\begin{lemma}[Lyapunov-Goodhart quantitative bridge]
\label{lem:2}
Under condition (C9) of Theorem~\ref{thm:deployment_safety}
(world-model parameterization regularity, comprising sub-conditions
BD bounded per-capability dimension count, CL coordinate-Lipschitz
parameterization, WB safety-relevant subspace with weight floor,
and TF truth floor) and under invariants $I_1$ and $I_4$ (which
together imply $I_2$ via
\citep[Phase~Redundancy]{lasser2026phase}'s contraction analysis),
\citep[Microfoundation]{lasser2026micro}'s relative sup-norm gap
$\epsnonres = \sup_{c \in P^{\mathrm{act}}} |P(c) - T(c)|/T(c)$ on
the operationally active subspace $P^{\mathrm{act}} = \{c \in P
\setminus \mathrm{ResS} : T(c) > 0\}$ satisfies:
\[
  \Lyap(\Wm_t) \;<\; \epsilon_{\mathrm{safe}}
  \quad\implies\quad
  \epsnonres \;<\; f(\epsilon_{\mathrm{safe}})
  \;:=\; \frac{L_{\max}}{T_{\min}}
  \sqrt{\frac{N_{\max} \cdot \epsilon_{\mathrm{safe}}}{w_{\min}}}.
\]
$f$ is intensive in $|P|$ over the deployment class $\mathcal{D}$:
each constant ($L_{\max}, T_{\min}, N_{\max}, w_{\min},
\epsilon_{\mathrm{safe}}$) depends on $\mathcal{D}$'s operational
parameters but not on $|P|$.
\end{lemma}

\begin{proof}[Proof outline; full proof in
Appendix~\ref{app:proof_lem2}]
By (C9.CL), the per-capability proxy-truth gap satisfies $|P(c) -
T(c)| \leq \sum_{k \in \mathrm{dim}(c)} L_{c,k} |\epsilon_k|$.
Apply weighted-dual-norm Cauchy-Schwarz with weights
$L_{c,k}/\sqrt{w_k}$ and $\sqrt{w_k} |\epsilon_k|$:
\[
  \sum_{k \in \mathrm{dim}(c)} L_{c,k} |\epsilon_k|
  \leq \sqrt{\sum_{k \in \mathrm{dim}(c)} L_{c,k}^2/w_k} \cdot
  \sqrt{\Lyap}.
\]
Bound the dual-norm factor by $L_{\max}^2 N_{\max}/w_{\min}$ via
(C9.BD) and (C9.WB), giving $|P(c) - T(c)| \leq L_{\max}
\sqrt{N_{\max} \Lyap / w_{\min}}$.  Apply (C9.TF) for the relative
sup-norm bound: $\epsnonres \leq (L_{\max}/T_{\min}) \sqrt{N_{\max}
\Lyap/w_{\min}}$.  Substitute $\Lyap < \epsilon_{\mathrm{safe}}$.
\end{proof}

\textbf{Connection to Theorem~\ref{thm:deployment_safety}'s Layer~1.}
Lemma~\ref{lem:2} converts the Lyapunov bound (Phase Redundancy
deliverable under invariants $I_1, I_4$) into the relative sup-norm
gap (Microfoundation Goodhart-bound input).  Theorem~1's Layer~1
proof invokes Lemma~\ref{lem:2} at Step~3, then composes via
gap decomposition (Step~4), Goodhart bound (Step~5), and Lemma~1's
generic-X intensivity packaging (Step~6) with $X =
f(\epsilon_{\mathrm{safe}})$.

\subsection{Lemma 3 (HHI-pressure surrogate adequacy)}
\label{sec:lem3}

We separate what the HHI argument actually proves (one-way
implication) from what the theorem requires (the reverse
implication, which is an empirical claim).

\begin{lemma}[HHI as concentration surrogate, forward implication]
\label{lem:3_forward}
Let $\mathcal{R}_{\mathrm{press}}$ denote the optimization-pressure
regime in which \citep[Microfoundation]{lasser2026micro}'s Conjecture~1 predicts gap exploitation.
Deployments with $\HHI > H^*$ exhibit trade-flow concentration
patterns that match the structural prerequisites of $\mathcal{R}_{\mathrm{press}}$:
\[
  \HHI > H^* \;\Longrightarrow\;
  \text{trade-flow concentration prerequisites of }
  \mathcal{R}_{\mathrm{press}} \text{ are met}.
\]
\end{lemma}

\begin{proof}
\citep[Need~Sufficiency]{lasser2026paper8b}'s HHI is the Schur-convex concentration index on trade-flow
events. By
\citep[Schur-convexity property]{lasser2026paper8b}, $\HHI > H^*$
implies that trade flow is concentrated on a small subset of
capability categories. Under stance S0, trade flow is the
operational signature of the agent's exercise pattern (\citep[Revealed~Sacrifice]{lasser2026paper8a}),
which is the operational signature of optimization target.
Therefore $\HHI > H^*$ implies optimization-target concentration,
which is predicate (a) of Conjecture~1's predicted-exploitation
regime.
\end{proof}

\paragraph{The reverse implication required by Theorem~\ref{thm:deployment_safety}.}
Lemma~\ref{lem:3_forward} establishes only a sufficient condition
for $\mathcal{R}_{\mathrm{press}}$. The deployment-safety theorem's
Layer~1 proof (\S\ref{sec:proof_layer1}, Step~4) requires the
\emph{contrapositive}: $\mathcal{R}_{\mathrm{press}}^c$ implies
$\HHI \leq H^*$, equivalently $\HHI < H^* \implies
\mathcal{R}_{\mathrm{press}}^c$. This direction does \emph{not}
follow from Lemma~\ref{lem:3_forward}. There may exist deployments
in $\mathcal{R}_{\mathrm{press}}$ with low HHI (e.g., concentrated
optimization pressure that does not manifest in concentrated trade
flow). For Theorem~1 to invoke $I_5$ as the operational invariant,
this gap must be closed.

We close it by an explicit assumption rather than by argument:

\begin{assumption}[HHI surrogate adequacy, SA1]
\label{ass:SA1}
For deployments under stance S0 with operational invariants $I_1,
\ldots, I_4$ in force, the HHI ceiling $\HHI < H^*$ is sufficient to
ensure the deployment lies outside Conjecture~1's predicted-
exploitation regime $\mathcal{R}_{\mathrm{press}}$:
\[
  \HHI < H^* \;\Longrightarrow\; \mathcal{R}_{\mathrm{press}}^c.
\]
\end{assumption}

\textbf{What SA1 buys vs.\ what proving it would buy.}
Under SA1, the deployment-safety theorem's Layer~1 proof binds via
$I_5$. The theorem becomes conditional on SA1 in the same sense it
is conditional on the eleven invariants: SA1 is a stated assumption,
not a derived consequence. SA1 is empirically testable
(comparing actual gap-exploitation behavior against HHI levels in
real deployments) and is the natural target for the
HHI--$\Delta r_K$ structural conjecture deferred to follow-up work
(\S\ref{sec:conjecture_op}).

\textbf{What we cannot do.}  Without SA1, the theorem cannot use
$I_5$ to bound $\mathcal{R}_{\mathrm{press}}$.  The alternative is
to weaken the theorem to bound only the regime characterized by
forward HHI implications; this loses the operational utility of
$I_5$ as a deployment-time check.  SA1 is the minimum assumption
required to retain operational utility while keeping the
implication direction explicit.

\paragraph{Note on the Conjecture-dependence audit (\S\ref{sec:conjecture_audit}).}
The Conjecture-dependence audit listed $I_5$ as the only
conjecture-dependent invariant. SA1 makes this dependence explicit.
The other ten invariants are theorem-proof conditions derived from
established source-paper machinery; $I_5$ is a deployment-regime
condition whose theorem-utility is conditional on SA1.

\subsection{Lemma 4 (SPRT detection lead-time tail bound)}
\label{sec:lem4}

\begin{lemma}[SPRT lead-time tail bound]
\label{lem:4}
Let $T_{\mathrm{detect}}$ be the SPRT detection time for a
violation producing post-clipping LLR drift $\delta > 0$ under
(C5.IID).  Then for $t > A/\delta$:
\[
  \Pr[T_{\mathrm{detect}} > t]
  \;\leq\; \exp\!\left(-\frac{2(t\delta - A)^2}{t \Rh^2}\right),
\]
where $A = \log((1-\beta)/\alpha)$ is Paper~5's SPRT threshold,
$\Rh = b - a = 2\Bclip$ is the Hoeffding range width of the
clipped LLR (called $\Rh$ in Lemma~\ref{lem:6}; equal to the
$R$ used in this appendix), and $\delta$ is the post-clipping drift
floor.  Asymptotically ($t\delta \gg A$), this simplifies to
$\exp(-2T\delta^2/\Rh^2) = \exp(-\kappa T \delta)$ with
$\kappa = 2\delta/\Rh^2$.
\end{lemma}

\begin{proof}[Proof outline; full proof in Appendix~\ref{app:proof_lem4}]
\citep[Exogenous~Verification]{lasser2026exo}'s SPRT machinery
(\citep[Definition of Behavioral Consistency Monitor]{lasser2026exo})
gives the running log-likelihood ratio $\Lambda_t$ with per-step
drift $\delta$ under the alternative.  Wald's identity gives the
mean bound $\mathbb{E}[T_{\mathrm{detect}}] \approx A/\delta$;
Hoeffding-Azuma's inequality applied to bounded clipped
LLR increments (under (C5.SPRT) sub-clauses) gives the tail bound
above.

\textbf{Note on $\Lambda_t < A$.}  In the proof, $\Lambda_t < A$ is
\emph{necessary} for $T_{\mathrm{detect}} > t$ (i.e., the LLR
process has not yet crossed the threshold), not sufficient; the
event $\{T_{\mathrm{detect}} > t\}$ requires that the threshold has
not been crossed at any earlier time.  The tail bound is
nonetheless valid because the necessary condition's probability
upper-bounds the joint event's probability.

\textbf{Caveat.} \citep[Phase~Redundancy]{lasser2026phase}'s
$\tau_{\mathrm{meta}} \gtrsim C/I_k$ is stated as a scaling law,
not a theorem-level bound; Lemma~\ref{lem:6}'s lead-time-before-
cascade composition uses (C11.CLK)'s deterministic floor $\Ncasc$
plus calibration probability $\bclk$ in lieu of treating
$\tau_{\mathrm{meta}}$ as a hard inequality.
\end{proof}

\subsection{The Lemma 5 family: substrate-targeting adversarial events}
\label{sec:lem5_family}

The five-part Lemma~5 family covers substrate-targeting adversarial
events specifically. The family structure mirrors the three-layer
deployment claim: Lemma~5c provides the static safe region,
Lemma~5b provides the channel-restricted detection class,
Lemma~5d composes detection with the lead-time guarantee,
Lemma~5a supplies the substrate floor that both layers depend on,
and Lemma~5e extends Lemma~5b's machinery to environment-side
observables.

The family is restricted to substrate-targeting adversaries
(Definition~\ref{def:joint_indep}); other shock classes
(capability-targeting, coalition-internal corruption, environment
manipulation outside witness coverage) require separate machinery
or are acknowledged residuals. \S\ref{sec:scope} of the
introduction stated this scope explicitly.

\subsection{Lemma 5a (Substrate floor)}
\label{sec:lem5a}

\renewcommand{\thelemma}{5a}
\begin{lemma}[Substrate floor]
\label{lem:5a}
Under invariant $I_6'$ ($\mindep \geq m^* \geq 3$) and condition
(C10) of Theorem~\ref{thm:deployment_safety} (per-pair cooperative-
novelty rate floor with superposition; comprising sub-clauses
(C10.CN) heterogeneous per-pair rate floors $\rho_{ij} \geq \rho_0
> 0$ on audited subset $\Saudit \subseteq \mathcal{S}$ of $m^*$
substrates, and (C10.SU) pairwise channel superposition: disjoint
attribution + non-rivalrous production + joint-deployment
intensities), the cross-substrate cooperative novelty rate satisfies:
\[
  \rext \;\geq\; r_*(m^*) \;:=\; \rho_0 \cdot \binom{m^*}{2} \;>\; 0.
\]
For $m^* = 3$, $r_*(3) = 3\rho_0 > 0$.

The floor $r_*(m^*)$ is intensive in $|P|$ over the deployment
class $\mathcal{D}$: $\rho_0$ depends on $\mathcal{D}$'s
operational parameters (per-pair calibrated rates) but not on
$|P|$.
\end{lemma}

\begin{proof}[Proof outline; full proof in
Appendix~\ref{app:proof_lem5a}]
By $I_6'$, $\mindep \geq m^*$ substrates with joint failure-
correlation independence; (C10) certifies an audited subset
$\Saudit$ of $m^*$ substrates.  Pairwise cooperative channels
$\mathcal{C}^{(s_i, s_j)}$ are defined with exact-two-substrate
causally-necessary support; auxiliary post-production participants
do not change classification.  Under (C10.CN), each pair satisfies
$\rext^{(s_i, s_j)} \geq \rho_0 > 0$.  Under (C10.SU), the
pairwise channels are disjoint event classes with non-rivalrous
production, so per-pair rates sum:
\[
  \rext \;\geq\; \rext^{\mathrm{pairwise}} \;\geq\; \binom{m^*}{2}
  \cdot \rho_0 \;=\; r_*(m^*).
\]
Higher-order cooperative contributions are nonnegative, so the
pairwise sum is a lower bound.
\end{proof}

\textbf{Connection to Theorem~\ref{thm:deployment_safety}'s
Layer~1.}  Lemma~5a's $\rext \geq r_*(m^*) > 0$ is invoked at
Step~1 of Lemma~5c's appendix proof to make the safe-region
inequality $\Delta r_K < \rext + (1-\gamma)(\Ddiv - \Delta_0)$
non-trivial.  Without (C10), the safe region could be vacuously
empty; with (C10), $\rext$ has a positive deployment-class
constant floor.

\textbf{(C10.SU) vs.\ joint failure-correlation independence ($I_6'$).}
$I_6'$ excludes shared adversarial \emph{shock} mechanisms: no
single shock spans multiple substrates.  (C10.SU) excludes shared
cooperative \emph{production} bottlenecks: producing cooperatives
in one channel does not consume resources that would otherwise
contribute to another.  These are structurally independent
conditions.  $I_6'$ alone does not imply additivity of cooperative
rates; (C10.SU) is required for that.

\subsection{Lemma 5b (Channel-restricted detection KL floor)}
\label{sec:lem5b}

\addtocounter{lemma}{-1}\renewcommand{\thelemma}{5b}
\begin{lemma}[Channel-restricted detection]
\label{lem:5b}
Let $\Aadv$ denote the subset of regime-(iii) strategies producing
at least one monitored four-channel deviation: reduced
cross-substrate cooperative event rate (Channel~1), degraded
attestation share (Channel~2), individuation/substrate-diversity
drift (Channel~3), or bundle/cooperative-edge decomposition shifts
(Channel~4). Then for any $q \in \Aadv$:
\[
  D_{\mathrm{KL}}(q \,\|\, p_0) \;\geq\; \deltaadv \;>\; 0,
\]
where $\deltaadv$ is computed from channel-specific
least-favorable alternatives.
\end{lemma}

\begin{proof}[Proof outline; full proof in Appendix~\ref{app:proof_lem5b}]
For each of the four channels, a least-favorable adversarial
distribution is constructed (Poisson for the cooperative-rate
channel, Bernoulli for attestation, multinomial for the two
concentration channels) and its KL divergence is computed against
the baseline.  Taking $\deltaadv = \min_c \delta_c$ over the four
channels gives a strictly positive lower bound on
$D_{\mathrm{KL}}(q \,\|\, p_0)$ for $q \in \Aadv$.

\textbf{The named gap (channel-orthogonal residual).}  Strategies
outside $\Aadv$ --- those producing no shift in any of the four
monitored channels --- have $D_{\mathrm{KL}}(q \,\|\, p_0) = 0$ on
the monitored observables and are outside the detection guarantee.
This is a real, named residual; \S\ref{sec:scope} of the
introduction acknowledged it and \S\ref{sec:discussion} returns to
it.
\end{proof}

\subsection{Lemma 5c (Minimax static tightening)}
\label{sec:lem5c}

We internalize the key definitions and the single-shock statement
in the body; the multi-shock extension appears as a Remark below.

\begin{definition}[Counterfactual shock-loss fraction]
\label{def:loss_fraction}
For a coalition state $G_K$ and substrate $s \in S$:
\[
  \alpha_s(G_K) \;=\; \frac{\volL(G_K) - \volL(G_K^{-s})}{\volL(G_K)},
\]
where $G_K^{-s}$ is the coalition's capability poset with all
capabilities (individual and cooperative) supported on $s$ removed.
The worst-case shock-loss fraction under adversarial targeting is
$\ell^{\max}(G_K) = \max_s \alpha_s(G_K)$.

The shock-loss profile $(\alpha_1, \ldots, \alpha_m)$ is a vector of
independent counterfactuals, not a partition; with cross-substrate
cooperative dependencies, $\sum_s \alpha_s$ may exceed 1, while each
$\alpha_s \in [0, 1]$.
\end{definition}

\begin{definition}[Cooperative-overlap regime]
\label{def:coop_overlap}
A deployment is in the \emph{cooperative-overlap regime} if its
substrate-cooperative structure is dominated by cross-substrate
cooperative capabilities (cooperatives whose production requires
participation from multiple distinct substrates) rather than by
redundant capabilities (capabilities replicated across substrates).
Operationally: the $\volL$-weighted sum of cross-substrate
cooperatives exceeds the $\volL$-weighted sum of redundant
capabilities. The cooperative-vs-redundancy audit
(\S\ref{sec:deployment_tooling}) verifies this regime.
\end{definition}

\addtocounter{lemma}{-1}\renewcommand{\thelemma}{5c}
\begin{lemma}[Minimax static tightening, single-shock cooperative-overlap regime]
\label{lem:5c}
Restrict scope to substrate-targeting adversarial events under the
single-shock model ($N = 1$ a.s.) and to deployments in the
cooperative-overlap regime
(Definition~\ref{def:coop_overlap}). Under invariants $I_6'$
($\mindep \geq m^* \geq 3$) and the balanced-loss condition:
\[
  \max_s \alpha_s(G_K^{\mathrm{div}}) \;\leq\; 1/m^* + \epsilon
  \qquad (\epsilon > 0 \text{ engineering tolerance}),
\]
\citep[Horizon~Aware]{lasser2026horizon}'s anti-monopolar conclusion holds in the static safe region:
\begin{equation}
\label{eq:safe_region_paper}
  \Delta r_K \;<\; \rext + (1 - \gamma)\bigl(\Ddiv - \Delta_0\bigr),
\end{equation}
where
\begin{equation}
\label{eq:ddiv_paper}
  \Ddiv \;=\; \volL_K \cdot (\lossDmax - \lossdivmax) \cdot
  \mathbb{E}[\gamma^{\Tadv}],
\end{equation}
$\lossDmax = \max_s \alpha_s(G_K^{D})$ is the dominator's
worst-case shock-loss fraction (approximately 1 under full
failure-correlated single-substrate domination), and
$\lossdivmax = \max_s \alpha_s(G_K^{\mathrm{div}})$ is the
diversity strategy's worst-case shock-loss fraction (bounded above
by $1/m^* + \epsilon$ under the balanced-loss condition).
\end{lemma}

\begin{proof}[Proof outline; full proof in
Appendix~\ref{app:proof_lem5c}]
The proof composes \citep[Horizon~Aware]{lasser2026horizon}'s strategy-dependent corollary with the
substrate-targeting shock model under the balanced-loss condition.

\emph{Step 1.} Add a substrate-targeting shock at random time
$\Tadv$ to \citep[Horizon~Aware]{lasser2026horizon}'s deterministic-growth model.

\emph{Step 2.} Compute the substrate-diversification value
advantage as the discounted expected difference in surviving
$\volL$:
$\Ddiv = \volL_K(\lossDmax - \lossdivmax) \mathbb{E}[\gamma^{\Tadv}]$.

\emph{Step 3.} Under $I_6'$ and the balanced-loss condition,
$\lossdivmax \leq 1/m^* + \epsilon$.

\emph{Step 4.} Add $(1-\gamma)\Ddiv$ to \citep[Horizon~Aware]{lasser2026horizon}'s strategy-dependent
corollary inequality. Diversity strictly dominates when the sum is
positive, yielding Equation~\ref{eq:safe_region_paper}.

\emph{Step 5 (cooperative-overlap regime conservativeness).}  In
the cooperative-overlap regime, sequential cooperative loss
overcounts when computed via original-counterfactual $\alpha_s$
(cooperatives are double-counted across the substrates that produce
them).  The equal-substrate analysis using the original $\alpha_s$
therefore underestimates diversity's surviving volume relative to
the actual sequential dynamics, making the safe-region bound a
\emph{lower} bound on the true advantage --- conservative for the
deployment claim.
\end{proof}

\begin{remark}[Multi-shock extension]
\label{rem:multishock}
For high-rate adversarial environments where multiple shocks
accumulate within the planning window, the single-shock $\Ddiv$
extends to:
\[
  \Ddiv^{\mathrm{multi}} \;=\; \mathbb{E}\!\left[\sum_{i=1}^{\infty}
  \gamma^{T_i} \bigl(\delta_i^D - \delta_i^{\mathrm{div}}\bigr)\right],
\]
where $\delta_i^{D}, \delta_i^{\mathrm{div}}$ are the per-shock
marginal $\volL$-losses for the domination and diversity strategies.
The single-shock case ($N = 1$ a.s.) reduces to
Equation~\ref{eq:ddiv_paper}.

Under independent substrate-targeting shocks at Poisson rate
$\lambda$, the multi-shock advantage admits a closed form
\[
  \Ddiv^{\mathrm{multi}} \;\geq\; \volL_K \left[\kappa -
  \frac{1}{\mindep} \cdot \frac{\kappa(1 - \kappa^{\mindep})}
  {1 - \kappa}\right]
\]
with $\kappa = \mathbb{E}[\gamma^{T_1}] = \lambda/(\lambda - \ln \gamma)$.
The closed form has a removable singularity at $\kappa = 1$;
$\Ddiv^{\mathrm{multi}}/\volL_K \to 0^+$ as $\kappa \to 1$. The
operational regime condition for non-trivial multi-shock advantage
is $\kappa < 1 - \delta$ for stated $\delta > 0$.

The multi-shock extension is used only for high-$\lambda$
deployments (adversarial environments with frequent attacks); the
single-shock form is the binding inequality for typical deployments.
The full multi-shock derivation is given in the multi-shock
paragraph of Appendix~\ref{app:proof_lem5c}.
\end{remark}

\subsection{Lemma 5d (Lead-time tail composition)}
\label{sec:lem5d}

\addtocounter{lemma}{-1}\renewcommand{\thelemma}{5d}
\begin{lemma}[Lead-time tail composition (union-class)]
\label{lem:5d}
Combining Lemmas~5b and~5e's KL floors with Lemma~4's tail bound,
the SPRT detection of any $q \in \Aadv \cup \Aadv^{\mathrm{env}}$
violation satisfies the union-class lead-time guarantee.  Under
(C5.SPRT) sub-clauses, define $\deltastar$ as the post-clipping
union-class drift floor satisfying $\deltastar \leq \min(\deltaadv,
\deltaadv^{\mathrm{env}})$.  Then by Lemma~\ref{lem:6}'s
Hoeffding-inversion derivation, the SPRT detection quantile
$\Tbeta(\beta, \deltastar)$ satisfies
\[
  \Pr[T_{\mathrm{detect}} > \Tbeta] \;\leq\; \beta.
\]
Composing with the metastable cascade lower bound
$T_{\mathrm{cascade}} \geq \tau_{\mathrm{meta}}$
(\citep[Phase~Redundancy]{lasser2026phase} Proposition on metastable
lifetime) and (C11.CLK)'s clock-comparability with calibrated
probabilities $\bclk, \bcal$:
\[
  \Pr[\Tbeta \leq \tau_{\mathrm{meta}}] \;\geq\; 1 - \bclk - \bcal,
\]
giving the lead-time-before-cascade guarantee.  Total Layer~2
failure probability is $\beta + \bclk + \bcal$.
\end{lemma}

\begin{proof}
The composition is direct, in three steps:

\emph{Step 1 (union-class drift floor).}  For $q \in \Aadv$,
Lemma~\ref{lem:5b} gives $D_{\mathrm{KL}}(q \,\|\, p_0) \geq
\deltaadv > 0$; for $q \in \Aadv^{\mathrm{env}}$,
Lemma~\ref{lem:5e} gives $D_{\mathrm{KL}}(q \,\|\, p_0^{\mathrm{env}})
\geq \deltaadv^{\mathrm{env}} > 0$.  Under (C5.SPRT) sub-clauses
(in particular (C5.HOEFF) clipping and (C5.IID) post-clipping
adapted increments), the clipped LLR process has conditional drift
$\geq \deltastar$ for any $q$ in the union class, with $\deltastar
\leq \min(\deltaadv, \deltaadv^{\mathrm{env}})$.  Clipping can
reduce drift; the inequality is uniform.

\emph{Step 2 (SPRT detection quantile from Lemma~6).}  By
Lemma~\ref{lem:6} (Substeps 5a--5e in
Appendix~\ref{app:proof_lem6}), the SPRT high-probability detection
quantile $\Tbeta = \max\{2A/\deltastar, 2\Rh^2 \log(1/\beta) /
\deltastar^2\}$ rigorously satisfies $\Pr[T_{\mathrm{detect}} >
\Tbeta] \leq \beta$ via Lemma~\ref{lem:4}'s Hoeffding form.

\emph{Step 3 (cascade-clock composition with (C11.CLK)).}
(C11.CLK) provides $\Pr[\Nev(\tau_{\mathrm{meta}}) \geq \Ncasc]
\geq 1 - \bclk$ with deterministic audit-time inequality $\Ncasc
\geq \Tbeta$ (certified case $\bcal = 0$; with calibration-
uncertainty $\bcal$ in the empirical case).  Chaining
$\Tbeta \leq \Ncasc \leq \Nev(\tau_{\mathrm{meta}})$ gives
$\Pr[\Tbeta \leq \tau_{\mathrm{meta}}] \geq 1 - \bclk - \bcal$.

\emph{Step 4 (union bound for total Layer 2 failure).}  By union
bound on (a) detection-tail $\{T_{\mathrm{detect}} > \Tbeta\}$
(probability $\leq \beta$) and (b) cascade-clock event
$\{\Tbeta > \tau_{\mathrm{meta}}\}$ (probability $\leq \bclk +
\bcal$):
\[
  \Pr[\text{Layer 2 fails}] \;\leq\; \beta + \bclk + \bcal.
\]

The composition exposes a clean operational regime: deployments
with long metastable cascade time and large post-clipping drift
floor have exponentially small detection-tail $\beta$; the
clock-comparability $\bclk$ and calibration $\bcal$ failures are
operator-tunable via Audit~7
(\S\ref{sec:tooling_substrate_audits}).  $\Box$
\end{proof}

\subsection{Lemma 5e (Environment-side witness extension)}
\label{sec:lem5e}

\addtocounter{lemma}{-1}\renewcommand{\thelemma}{5e}
\begin{lemma}[Environment-side witness extension]
\label{lem:5e}
Under invariant $I_8$ and theorem-level condition (C12.ENV-WIT)
with its six sub-clauses, Lemma~\ref{lem:5b}'s channel-restricted
detection class extends to environment-side observables.
Specifically, for any $q \in \Aadv^{\mathrm{env}}$ producing a
threshold-exceeding shift in some $v \in \Vem$:
\[
  D_{\mathrm{KL}}(q \,\|\, p_0^{\mathrm{env}}) \;\geq\;
  \deltaadv^{\mathrm{env}} := \min_{v \in \Vem}
  \delta_v^{\mathrm{env}} \;>\; 0,
\]
where $\delta_v^{\mathrm{env}}$ is the per-variable least-favorable
KL floor (Poisson for rate variables, Bernoulli for indicator
variables, with support conventions from (C5.SUPP)).
\end{lemma}

\begin{proof}[Proof outline; full proof in Appendix~\ref{app:proof_lem5e}]
The proof has four parts:
(a)~the symmetric construction inheriting Paper~5's
substrate-exclusivity machinery under the (C12.ENV-WIT) trust
model (witness substrate outside adversary write-access and
failure-correlation domains);
(b)~per-variable KL floor derivations for the canonical four
variables $V_1, V_2, V_3, V_4$ (Poisson cooperative rate $\rext$,
Bernoulli substrate-distinctness, Poisson adversarial-arrival
rate $\lambda$, Bernoulli trusted-setup status), each strictly
positive under (C5.SUPP) support conventions;
(c)~the data-processing inequality applied to the fixed
pre-deployment witness-recording map gives
$D_{\mathrm{KL}}(q \,\|\, p_0^{\mathrm{env}}) \geq
D_{\mathrm{KL}}(q_{v_i} \,\|\, p_{0, v_i}) \geq
\delta_i^{\mathrm{env}}$;
(d)~strict-positivity of $\deltaadv^{\mathrm{env}}$ as the minimum
over a finite (per (C12.PUB)) set of strictly positive
quantities.

\textbf{Composition with Lemma 6.}  The post-clipping union-class
drift floor $\deltastar$ used in Lemma~\ref{lem:6}'s $\Tbeta$
satisfies $\deltastar \leq \min(\deltaadv,
\deltaadv^{\mathrm{env}})$, with the inequality accounting for any
drift reduction from (C5.HOEFF) clipping.  Lemma~5e supplies
$\deltaadv^{\mathrm{env}}$ for the union; Lemma~5b supplies
$\deltaadv$.

\textbf{The named residual (R2), refined.}  R2 covers manipulations
satisfying \emph{both} (i)~target only exogenous variables outside
$\Vem$, \emph{and} (ii)~produce no threshold-exceeding shift in any
monitored $v$ (direct or causal).  Manipulations of unmonitored
variables that causally shift a monitored $v$ are detected through
that monitored projection (not R2).
\end{proof}

\subsection{Lemma 6 ($h_{\mathrm{detect}}$ intensivity)}
\label{sec:lem6}

\renewcommand{\thelemma}{\arabic{lemma}}
\begin{lemma}[$h_{\mathrm{detect}}$ intensivity]
\label{lem:6}
Under (C5) continuous SPRT monitoring, (C5.SPRT) sub-clauses
(C5.HOEFF/MULT/IID/SUPP), (C11) bounded gap-growth rate, and
(C11.CLK) clock comparability, combined with Lemma~\ref{lem:4}
and Lemma~\ref{lem:1}: the SPRT high-probability detection quantile
\[
  \Tbeta(\beta, \deltastar) \;:=\;
  \max\!\left\{\frac{2A}{\deltastar},\;
  \frac{2 \Rh^2 \log(1/\beta)}{\deltastar^2}\right\}
\]
satisfies $\Pr[T_{\mathrm{detect}} > \Tbeta] \leq \beta$, and on
the joint event $\{T_{\mathrm{detect}} \leq \Tbeta\} \cap \{\Tbeta
\leq \Nev(\tau_{\mathrm{meta}})\}$:
\[
  h_{\mathrm{detect}}(\theta) := \sup_{s \in [t_0, t_0 +
  T_{\mathrm{detect}}]} \epsgap(s) \;\leq\;
  h_{\mathrm{static}}(\theta) + \rhogap(\deltastar) \cdot \Tbeta.
\]
$h_{\mathrm{detect}}(\theta)$ is intensive in $|P|$ over
$\mathcal{D}$.  Total Layer~2 failure probability is $\beta + \bclk
+ \bcal$ (with $\bcal = 0$ in the certified-conservative case).

Here $A = \log((1-\beta)/\alpha)$ is Paper~5's SPRT threshold,
$\Rh = 2\Bclip$ is the (C5.HOEFF) Hoeffding range width, and
$\deltastar$ is the post-clipping union-class drift floor
((C5.IID), Lemmas~\ref{lem:5b}, \ref{lem:5e}).
\end{lemma}

\begin{proof}[Proof outline; full proof in Appendix~\ref{app:proof_lem6}]
The proof has two roles for cascade time, kept separate.  $\Tbeta$
is the integration horizon for $\rhogap$, derived directly from
Lemma~\ref{lem:4}'s Hoeffding form via a rigorous (non-asymptotic)
chain: $T \geq 2A/\deltastar$ implies $T\deltastar - A \geq
T\deltastar/2$, and squaring gives $(T\deltastar - A)^2 \geq
T^2\deltastar^2/4$, so the Hoeffding exponent $\geq T\deltastar^2/
(2\Rh^2)$.  Solving for exponent $\geq \log(1/\beta)$ gives
$T \geq 2\Rh^2\log(1/\beta)/\deltastar^2$.  The $\max$-form
$\Tbeta = \max\{2A/\deltastar, 2\Rh^2\log(1/\beta)/\deltastar^2\}$
covers both regimes.

For the supremum bound: under the boundary convention $t_0 := t_0^-$
(the last safe SPRT step), Lemma~\ref{lem:1} gives $\epsgap(t_0)
\leq h_{\mathrm{static}}(\theta)$.  (C11) bounds the per-step
post-clipping increment by $\rhogap(\deltastar)$, giving $\epsgap
\leq h_{\mathrm{static}} + n \cdot \rhogap$ for $n$ steps.

For the cascade-clock event: (C11.CLK) provides
$\Pr[\Nev(\tau_{\mathrm{meta}}) \geq \Ncasc] \geq 1 - \bclk$ with
$\Ncasc \geq \Tbeta$ as a hard audit constraint (in the certified
case; with calibration-uncertainty $\bcal$ in the empirical case).
The chain $\Tbeta \leq \Ncasc \leq \Nev(\tau_{\mathrm{meta}})$ gives
the lead-time-before-cascade guarantee with probability $\geq 1 -
\bclk - \bcal$.
\end{proof}

\textbf{Constant taxonomy.}  $\Tbeta$ depends on $\alpha$ (Type-I
rate, deployment-class via (C5.MULT) familywise allocation),
$\beta$ (detection-tail parameter, monitor-design), $\Bclip$ (clip
radius, (C5.HOEFF), deployment-class), $\Rh = 2\Bclip$ (Hoeffding
range, Lemma~\ref{lem:4}'s $R = b - a$), and $\deltastar$
(post-clipping drift floor, (C5.IID), Lemmas~5b/5e).  Cascade-clock
event uses $\Ncasc, \bclk, \bcal$ from (C11.CLK).  All
deployment-class.

\textbf{Sign-direction discipline.}  All operator certificates point
conservatively: $\Bclip$ \emph{upper}, $\deltastar$ \emph{lower},
$\Ncasc$ \emph{lower}, with calibration-uncertainty $\bcal$ in the
empirical case.

\subsection{Lemma family composition}
\label{sec:lem_composition}

The five-part Lemma~5 family composes as follows in the
deployment-safety theorem (\S\ref{sec:main_theorem}):

\begin{itemize}
\item \textbf{Static layer}: Lemma~5c supplies the safe region
  inequality. Inside the safe region, no substrate-targeting
  adversarial event (within the cooperative-overlap regime)
  produces a violation of the anti-monopolar conclusion.
\item \textbf{Detection layer}: Lemmas~5b, 5d, and~6 compose to give
  the channel-restricted detection class with tail-bounded lead time.
  Substrate-targeting evasions of the safe region within $\Aadv \cup
  \Aadv^{\mathrm{env}}$ are detected with
  $\Pr[T_{\mathrm{detect}} > \Tbeta] \leq \beta$, and the
  cascade-clock event $\{\Tbeta \leq \Nev(\tau_{\mathrm{meta}})\}$
  holds with probability $\geq 1 - \bclk - \bcal$.  Total Layer~2
  failure $\leq \beta + \bclk + \bcal$.
\item \textbf{Substrate floor}: Lemma~5a supplies the foundation
  for both layers ($\rext$ bounded below by $r_*(m^*) > 0$ given
  $\mindep \geq m^*$).
\item \textbf{Environment-side extension}: Lemma~5e extends the
  detection class to environment-side observables under invariant
  $I_8$.
\item \textbf{Acknowledged residuals}: strategies outside
  $\Aadv \cup \Aadv^{\mathrm{env}}$ (channel-orthogonal,
  environment-witness-orthogonal) and deployments outside the
  cooperative-overlap regime (Lemma~5c-prime open work) fall
  outside the guarantee.
\end{itemize}

The composition with Lemmas~1--4 supplies the bridge to \citep[Microfoundation]{lasser2026micro}'s
static Goodhart bound and the operational lead-time guarantee.
\S\ref{sec:main_theorem} states the full deployment-safety
theorem and traces the composition of all nine lemmas.
